% !TEX program = xelatex
\documentclass[UTF8, a4paper, 12pt]{ctexart}
\usepackage{geometry}
\usepackage{amsmath, amssymb}
\usepackage{booktabs}
\usepackage{enumitem}
\usepackage{setspace}
\usepackage{hyperref}
\usepackage[spaces]{xurl}
\geometry{left=2.5cm, right=2.5cm, top=2.5cm, bottom=2.5cm}
\hypersetup{colorlinks=true, linkcolor=blue, urlcolor=blue}
\onehalfspacing

% 行内路径与变量名：略小于正文、可在 '_' '/' 等处断行
\makeatletter
\g@addto@macro{\UrlBreaks}{\do\_}
\makeatother
\Urlmuskip=0mu plus 2mu
\newcommand{\id}[1]{{\footnotesize\nolinkurl{#1}}}

\title{程序化交易监管与股票市场微观结构\\实证研究中期报告（扩展稿）}
\author{（姓名 / 学号）}
\date{\today}

\begin{document}
\maketitle

\begin{abstract}
\noindent
程序化交易在 A 股成交中占比持续上升，监管规则在信息披露、报告义务与异常交易监控等方面对市场主体形成约束。本论文将问题聚焦于：在准自然实验框架下，二零二四年十月八日《证券市场程序化交易管理规定（试行）》施行前后，股票层面市场微观结构指标是否发生系统性变化；以及上述变化是否随政策前可观测的“程序化相关暴露”（以订单流不平衡等高频代理刻画）呈连续型异质性。中期阶段已完成从逐笔成交与限价订单簿到日频面板的可复现流水线，构造有效价差、实现波动率、Amihud 非流动性及政策前连续暴露，并在股票与交易日双向固定效应下完成主回归与多因变量、多暴露的网格稳健性程序。全样本下主交互项系数在双向聚类标准误下呈约百分之十水平的边际显著，推断对聚类设定敏感。报告在后文分节说明研究问题的逻辑分解、全文章节架构、识别策略与数据工程细节、初步结果及局限，并给出下一阶段在事件研究、协变量扩展、主因变量遴选披露与论文定稿方面的工作计划。
\end{abstract}

\tableofcontents
\newpage

\section{问题提出与研究意义}

\subsection{现实背景}

近年来，算法与程序化交易在境内证券市场的参与度上升，价格发现与流动性供给的效率收益与短期波动、公平性之间的张力受到监管者与学术界共同关注。中国证监会于二零二四年五月十一日发布《证券市场程序化交易管理规定（试行）》，并自二零二四年十月八日起施行；此后沪深北交易所于二零二五年七月七日发布《程序化交易管理实施细则》，对高频认定、监测指标等作进一步细化。制度上存在至少三个可区分的时点：发布（预期形成）、施行（规则生效与合规边界变化）、以及交易所细则落地（执行层面的技术细化）。若不加区分地将三者混为单一“政策虚拟变量”，容易混淆预期效应与合规约束效应，也会在事件研究中产生动态系数难以解释的问题。因此，论文在叙事与实证上须明确：主识别锚点与哪一时点对齐，其余时点如何进入稳健性或机制讨论。

\subsection{经验问题的分解}

本研究的核心经验问题可分解为三个层次。第一，平均而言，在控制股票个体效应与日历日共同冲击后，规章施行后全市场或样本内股票的微观结构指标（如交易成本代理、短期波动、非流动性）是否相对政策前发生显著变化。第二，上述变化是否在不同股票之间呈系统性异质性；特别是，事前在订单流形态上更接近“高强度程序化相关交易模式”的股票，是否在政策后呈现不同的轨迹。第三，在异质性存在的前提下，能否在平行趋势与无混淆等假设下赋予交互项系数以因果解释，并讨论机制（流动性补偿、报价竞争、监管约束的异质执行等）与政策含义的边界。

第一问对应传统“全样本政策前后对比”在时间固定效应下的吸收；第二问要求引入事前测度的连续暴露并与政策后虚拟变量交互，从而形成连续型双重差分式识别；第三问涉及计量文献中关于双向固定效应与异质性处理效应下系数含义的讨论，须在正文中与 de Chaisemartin 与 d'Haultfoeuille、Callaway 与 Sant'Anna、Sun 与 Abraham 等线索对接，诚实说明本文系数为特定加权平均处理效应框架下的概括，而非简单等价于“对每只股票的处理效应平均”。

\subsection{研究意义}

在学术层面，本文将境内程序化监管“组合拳”置于微观结构文献中市场质量与订单流信息的分析传统下，尝试在可得的逐笔与簿记数据约束下，把政策评估从离散分组（如指数成分）推进到事前连续暴露与政策交互的规格，并与多期 DID、事件研究方法的最新讨论对话。在政策层面，本文对“监管是否改善交易成本与稳定性”提供可复核的量化证据，并强调异质性结论对精准监管与差异化监测的潜在含义。在工程层面，本文坚持“高频计算一次物化、低频回归反复读取”的原则，使学位论文的实证部分具备可复现目录结构与命令行接口，降低合作者与答辩委员对“黑箱数据处理”的疑虑。

\section{制度背景与识别上的主要难点}

\subsection{与主回归对齐的时点选择}

代码与当前数据规格将主回归中的 \(\mathrm{Post}_t\) 对齐于规章施行日（二零二四年十月八日），而非发布日（二零二四年五月十一日）。理由包括：施行日是市场主体合规义务与报告边界发生实际变化的法律节点；发布日至施行日之间可能存在学习与预期交易，若主 \(\mathrm{Post}\) 取发布日，估计量混合了预期与生效后行为，解释负担更重。发布日更适合单独做事件研究窗口或作为安慰剂、预期检验的时间锚点。交易所细则（二零二五年七月七日）可作为扩展分析中的第二次冲击或机制变量来源，不宜与主施行日在同一张主表中不加区分地堆叠，以免交错采纳下 TWFE 系数解释进一步复杂化。

\subsection{事前暴露的内生性与测量误差}

事前连续暴露 \(\tilde{E}_i\) 由政策前窗口内基于逐笔成交计算的订单流不平衡等指标经时间平均、截面 winsor 与标准化得到。其经济学含义是“政策前该股成交流在买卖方向上的相对强度”，与“程序化交易强度”在观测上相关，但并非一一映射：部分主观交易、机构调仓亦可能呈现极端 OIB。因此，\(\tilde{E}_i\) 应被理解为程序化相关暴露的高频代理，而非监管报备意义上的真实程序化强度。测量误差与概念错位会使交互项系数向零偏误，需在局限部分讨论。另需警惕的是，若政策前窗口内暴露与不可观测的公司基本面或交易制度偏好相关，且这些因素同时影响政策后的结果，则即使加入股票固定效应，交互项识别仍可能受到混淆；缓解思路包括加入时变协变量、行业乘时间趋势、以及利用离散指数分组作对照 DID 检验结论是否由风格因素驱动。

\subsection{面板结构与误差相关}

数据为股票—交易日不平衡面板，同一股票多日观测误差通常序列相关，同一交易日多只股票受共同宏观与情绪冲击，误差横截面相关。忽略双重聚类往往导致标准误过小。中期全样本回归显示，仅按交易日聚类与双向聚类下同一系数点估计不变而显著性差异较大，说明推断设定对结论表述影响直接。正式论文应以双向聚类或理论上可辩护的更保守设定为基准报告，并并列给出单维聚类作对照，而非反选较窄标准误。

\section{论文整体架构}

全文拟按学位论文常规八章结构组织，各章职责如下。

第一章为引言，交代研究背景、制度时间线、核心经验发现梗概与边际贡献，并说明与开题报告的关系及中期后的修订要点。第二章为文献综述，分述微观结构中交易成本与订单流信息、算法交易与市场质量、监管评估的准自然实验方法，并单列一小节讨论“事前连续暴露与政策虚拟变量交互加双向固定效应”在方法与应用文献中的位置，以避免审稿人认为该规格为 ad hoc。第三章为制度背景与政策梳理，按时间顺序说明《规定》与交易所细则的条款要点，并明确本文主识别与哪些条款的经济机制相对应。第四章为数据与变量，给出逐笔与 LOB 的数据来源、清洗规则、样本区间（默认自二零二二年一月起至可得数据最后完整交易日止）、以及有效价差、实现波动率、Amihud、OIB 暴露等变量的操作化定义，附以与 Hendershott 等、Andersen 等经典文献的对照表。第五章为实证设计，写出基准 TWFE 方程、聚类方式、样本剔除规则、替代规格（事件研究、离散分组 DID、联合控制多种暴露等），并预先承诺主被解释变量从候选池中遴选的规则，以降低数据挖掘嫌疑。第六章报告主结果与稳健性，含动态系数图、替代因变量、替代暴露、子样本（如指数成分）与安慰剂。第七章讨论机制与经济含义，并审慎联系中小投资者保护叙事。第八章总结、政策建议与研究局限。

机器学习反事实或反事实构造若保留在开题中，则作为第六章或第七章的扩展小节，须明确与主 DID 的主辅关系及外生特征集约束，避免正文出现多种模型并列却无主次。中期阶段该部分尚未实现，下文在“下一步计划”中单独列出。

\section{识别策略与变量设计（与当前工程一致）}

\subsection{基准回归式}

基准模型为
\begin{equation}
y_{it} = \beta\, (\mathrm{Post}_t \times \tilde{E}_i) + \alpha_i + \gamma_t + \varepsilon_{it},
\end{equation}
其中 \(i\) 为股票（与交易所组合成的面板单元），\(t\) 为交易日；\(\alpha_i\) 与 \(\gamma_t\) 分别为股票与时间固定效应；\(\mathrm{Post}_t\) 在施行日及之后取 1；\(\tilde{E}_i\) 为政策前估计并截面标准化的 OIB 暴露，在数据列中记为 \id{E_oib_z}；交互项在数据列中记为 \id{post_x_E_oib_z}。被解释变量 \(y_{it}\) 的中期基准为成交量加权有效价差 \id{vw_es_bps}，候选池还包括 \id{rv_mid}、\id{amihud} 及成交活跃度类变量，由 \id{did-regress-grid} 批量估计以便稳健性表整理。

\subsection{离散 DID 的定位}

中证 1000 相对沪深 300 的离散双重差分可作为稳健性：检验连续暴露结论是否完全由市值—风格分层驱动。离散规格不作为主表核心估计量，以避免与连续规格争夺“主叙事”，但在答辩材料中准备一张对照表有助于回应偏好传统 TWFE 的读者。

\subsection{有效价差与暴露的计算要点（概念层）}

有效价差在实现上采用逐笔成交价与成交前一刻可得的最佳买卖报价中间价之偏离，按成交量加权聚合至日频；暴露采用政策前窗口内日度 OIB 的时间均值，再 winsor 与 z 化。二者均依赖 \id{trans_fea} 与 LOB 在同一交易日、同一交易所内的对齐质量。深圳委托文件路径在配置所给日期前后存在命名规则切换，工程上已在 \id{raw_paths} 中集中管理；若后续引入撤单率等扩展暴露，须避免与成交表路径混淆。

\subsection{与中小投资者保护叙事的衔接方式}

开题报告将政策评价与中小投资者保护相关联。中期阶段在实证上采取的做法是：先用微观结构文献中可辩护的市场质量代理（尤其是有效价差所映射的交易成本）建立主结果，再在讨论章将“价差收窄或波动变化”与订单执行成本、短期价格冲击等中间概念挂钩，从而间接涉及投资者福利；避免在系数刚取得边际显著时即作过强的福利论断。若后续引入订单流 toxicity 或逆向选择分量，机制链条可进一步收紧，但须与当前数据算力约束相权衡。

\section{文献与方法论文献线索（中期整理）}

微观结构实证方面，有效价差与买卖报价中间价的关系、以及成交量加权在日频聚合中的惯例，可追溯至 Hendershott、Jones 与 Menkveld 等关于算法交易与市场质量的工作；已实现波动率与日内收益波动的关系可参考 Andersen、Bollerslev 等与高频波动测度相关的研究。监管与事件研究方面，境内程序化交易相关实证仍相对分散，英文文献中关于算法交易报告义务、断路与熔断等规则的评估可作方法对照。

计量方法论方面，双向固定效应下异质性处理效应与交互项解释需谨慎，可引用 de Chaisemartin 与 d'Haultfoeuille 等；多期 DID 与交错采纳可引用 Callaway 与 Sant'Anna 等；事件研究中动态系数在异质性下的偏误讨论可引用 Sun 与 Abraham 等；平行趋势敏感性可引用 Rambachan 与 Roth 等。正文须在贡献段说明：本文主规格为施行日虚拟变量与事前连续暴露交互加 TWFE，与上述文献中的某一具体定理不必一一对应，但推断与解释上遵循同类警示。

\section{数据与工程进展}

\subsection{物化目录与版本}

派生数据根目录由 \id{config.py} 中 \id{DATA_ROOT} 与 \id{SCHEMA_VERSION} 决定，当前为 \path{.../prog_trade_reg/v1/}。指标定义变更时递增版本号，避免与旧结果混用。主要子目录包括：\id{daily/outcomes/}（按日 Parquet，两市合并，以 \id{exchange} 区分）、\id{meta/}（政策前暴露与构建清单）、\id{panel/}（DID 长表）。构建过程追加写入 \id{meta/build_manifest.jsonl}，便于审计跳过日期与失败原因。

\subsection{流水线命令与依赖顺序}

在配置好原始路径与交易日历 \id{trade_days_dict.pkl} 后，推荐顺序为：运行 \id{exposure} 生成 \id{meta/exposure_pre_policy.parquet}；运行 \id{outcomes-range} 或 \id{build-all} 批量生成 \id{daily/outcomes/*.parquet}；运行 \id{panel-did} 合并为 \id{panel/panel_did_long.parquet} 与 \id{panel_did_long.csv}；最后运行 \id{did-regress} 或 \id{did-regress-grid}。单日调试可用 \id{outcomes YYYYMMDD}。中期检查确认：暴露窗口结束日严格早于施行日；合并键为 \((\text{exchange}, \text{code})\)；\id{post} 由字符串交易日与施行日比较得到。须注意，批量任务对缺失或损坏的逐笔文件在默认模式下跳过该日，全样本覆盖率需在第四章以表格报告。

\subsection{回归软件与数值细节}

回归使用 Python 包 \id{linearmodels} 的 \id{PanelOLS}，公式中含实体与时间效应；大样本下启用低内存估计路径；对解释变量矩阵做秩检查，并在调试子样本时采用 bookend 或随机抽样，避免按日期排序文件仅取前若干行导致 \id{post_x_E_oib_z} 全为零进而秩亏。

\subsection{日度因变量的计算流程（与代码模块对齐）}

对每个交易日，程序读取合并后的逐笔 Feather \id{trans_fea}，按交易所后缀解析六位证券代码与整数内码，过滤成交价与成交量为正。对每一交易所分别读取当日 LOB Parquet，在 \((\text{code}, \text{time})\) 上排序并形成快照序列。逐笔成交以时间为键、按向后 as-of 规则匹配同代码下最近一条不晚于成交时刻的 LOB 记录，在买卖一价有效时计算中间价，进而得到逐笔有效价差（以基点度量），再按股票聚合为 \id{es_bps_num} 与 \id{es_bps_den}。同一 LOB 序列上可并行计算对数中间价差分的平方和，得到 \id{rv_mid}；结合美元成交额构造 Amihud 型日度指标 \id{amihud}。输出写入单日 Parquet，主键为 \((\text{exchange}, \text{code}, \text{trade\_date})\)。

\subsection{政策前暴露的聚合流程}

暴露模块对政策前区间内每个交易日读取 \id{trans_fea} 一次，在股票—交易所层面汇总主动买与主动卖成交量，计算日度 OIB，并可扩展成交笔数与总成交量。随后在时间维对每只股票求均值，在截面做分位数 winsor，再对 winsor 后的均值做全截面标准化，得到 \id{E_oib_z} 等列。该文件在行维度上为每只股票一行，与 outcome 日频表通过 \((\text{exchange}, \text{code})\) 左连接；连接后暴露列对同一股票在所有日历日上重复出现，与 \(\mathrm{Post}_t\) 相乘后在政策前恒为零，这是连续交互项识别结构的一部分，而非数据错误。

\subsection{长表合并与质量控制}

\id{panel-did} 扫描 \id{daily/outcomes/} 下全部 Parquet 并纵向拼接，与暴露表连接后生成 \id{post}、\id{policy_date_main}、\id{post_x_E_oib_z} 等列，并计算 \id{vw_es_bps}。质量控制要点包括：施行日与暴露窗口无重叠；两市代码不混用；缺失 outcome 或缺失暴露的股票日在回归中随 listwise deletion 剔除，须在正文报告剔除比例。CSV 导出使用竖线分隔以便 Stata 导入；大样本下优先使用 Parquet 做内存内估计以减少类型转换开销。

\section{初步实证结果}

基于已构建面板（全样本约数百万股票—日观测、股票维度千余只、交易日维度九百余日），\id{post_x_E_oib_z} 对 \id{vw_es_bps} 的系数在仅按交易日聚类时在常规水平上显著为负；在股票与交易日双向聚类下，点估计不变而标准误增大，显著性减弱至约百分之十边际水平，在百分之五水平上不再显著。非稳健 F 统计量与聚类稳健 F 统计量差异显著，说明若忽视误差相关结构会严重高估显著性；正文须并列报告稳健标准误与相应检验统计量。

经济含义上，负号与“事前 OIB 暴露较高的股票在政策后有效价差相对收窄”的方向一致。Within \(R^2\) 在双向固定效应吸收大量变异后通常极低，这不否定交互项系数的经济关注，但提示单一交互项对日内残差的解释比例有限，讨论经济显著性时应结合系数量级与 \id{vw_es_bps} 的尺度（基点）进行换算表述。

子样本或按文件首段截取曾得到更显著的系数，原因在于政策前期交互项在子样本内退化为常数或变异不足会导致调试规格失真；正式结论必须以全样本与预先承诺的聚类方式为基准。网格回归命令 \id{did-regress-grid} 可在同一面板读入下批量估计多种因变量与多种 \id{post_x_E_*} 交互项，中期可用于附录稳健性表的初稿生成，但最终入表结果须与正文叙事一致并避免事后挑选。

在双向固定效应框架下，交互项系数反映的是在吸收个股与时间均值之后，暴露与政策后共同变化所对应的相关结构；其因果解释仍依赖平行趋势与无混淆等识别假设。正文应引用异质性 TWFE 相关讨论，并避免将单一 p 值大小作为选择主因变量或聚类方式的唯一依据。

\section{当前局限}

其一，事件研究动态路径与平行趋势图尚未固化到与主规格同一套代码输出中，难以在答辩时展示系数随相对施行日的事件时间轨迹。其二，公司特征、行业、指数成分等协变量尚未系统并入主回归，异质性分析仍以连续暴露为主，对“风格混淆”的反驳仍不充分。其三，主被解释变量虽有多候选，但须在正文披露遴选规则与事前约定，目前文字尚未写死。其四，原始数据依赖集群路径，复现文档与匿名审稿材料需准备脱敏说明或公开子样本。其五，若学位论文保留机器学习反事实章节，与主 DID 的接口、外生特征集与训练窗尚未在代码层闭合。

\section{下一步工作计划（约八周毕业冲刺）}

前文“三至四个月”“中期再写某章”的安排与仅剩约两个月左右的答辩准备期不匹配。本节以答辩可完成为硬约束重写：先锁定最小可答辩闭环，再列争取完成项与建议暂缓项，避免在机器学习反事实、多软件复现、大而全文献扩写上分散精力。

\subsection{时间约束下的原则}

学位论文在八周量级内应以“一章一章写成可提交稿”为节奏，而不是再铺新的工程模块。实证主链条（数据—主回归—稳健性—事件研究图）优先于一切扩展；凡不直接出现在答辩 PPT 与纸质稿主文的内容，默认进附录或删节。与导师确认学院截稿、查重与送审日期，倒推两周留全文统稿与格式、一周留预答辩后修补。

\subsection{必须完成项（答辩硬门槛）}

\begin{enumerate}[leftmargin=2em]
  \item 以施行日为锚的事件研究或动态系数图至少一幅（可与主 TWFE 同一数据与样本），用于回应平行趋势与动态效应质疑；若代码尚未封装，优先用 Stata 或 R 在导出的面板上快速出图，再回到 Python 固化，避免卡在工程洁癖上。
  \item 主表定稿：双向聚类为主、按交易日聚类为对照；同一张表写清样本量、股票数、交易日数、剔除规则。
  \item 稳健性表最小集：全样本跑 \id{did-regress-grid} 或等价脚本，至少保留替代因变量（如 \id{rv_mid}、\id{amihud}）与至少一种替代暴露交互项各一列；指数或市值离散 DID 若时间紧可做一张小表或脚注级对照。
  \item 方法段事先写定：主因变量从候选池（有效价差、实现波动率、Amihud）中择一的规则落笔，避免事后挑选印象。
  \item 第三至六章成文：制度背景、数据与变量、实证设计、结果须叙事连贯；文献综述以够用、可引用、与主规格对话为度，不必追求穷尽。
  \item 摘要、结论、参考文献格式与学院模板一次对齐；替换本报告作者占位；准备查重与盲审所需匿名说明（数据路径可写“见复现包配置”）。
\end{enumerate}

\subsection{争取完成项（有时间再做）}

在必须项已可答辩的前提下：相对发布日的事件研究对照图；行业乘时间趋势或少量时变协变量；截距项在 TWFE 下报告方式的脚注；更完整的覆盖率与缺失表。Stata 或 R 与 Python 的系数交叉验证仅在与导师明确要求或盲审意见需要时投入，不作为八周内默认任务。

\subsection{建议暂缓或降级项（避免拖累毕业节点）}

\begin{enumerate}[leftmargin=2em]
  \item 开题中的机器学习反事实若尚未有闭合代码与稳定结果，建议在毕业稿中将该章整章降级为“未来研究”或附录极短述评，避免与主 DID 双线并行导致二者皆薄。
  \item 二零二五年七月细则作为第二次冲击的完整交错 DID，若样本期与叙事尚未理清，可不在主文展开，必要时一句稳健性交代即可。
  \item 大规模扩展暴露（撤单率、逐笔委托全链路）若未在现有流水线中跑通，不纳入主文，以免引入新的数据争议窗口。
\end{enumerate}

\subsection{建议周序（可按学院日历微调）}

第 1--2 周：事件研究或动态图出稿；主表与稳健性表数值冻结；与导师确认“主因变量择一”表述与答辩底线。第 3--4 周：第三至六章初稿串起来，图表编号与正文互引一致。第 5--6 周：机制与结论章写成可讨论稿；文献综述补齐与主规格相关的计量与应用引用；全文第一次统稿与查重。第 7 周：预答辩或院内模拟，按意见只改论证链路与格式，少改数据结构。第 8 周：定稿、打印、答辩 PPT 与备份数据与命令记录。

\subsection{答辩材料最小清单}

政策时间轴一页；主回归表（双向聚类为主）；事件时间图至少一幅；稳健性表一页；数据覆盖或缺失说明一段。仓库 \id{README.md} 与正文符号一致即可，不必在答辩前追求“第二套软件全复现”。

\subsection{风险与预案（保持简短）}

双向聚类下主系数若仍弱显著：以动态图与离散 DID 对照支撑“政策时点周边存在结构变化”的叙述，并在局限中诚实讨论推断与测量误差。符号异常时先查施行日与样本窗，再改解释，不改数据规则。

\subsection{复现与环境}

在目标环境中执行 \id{pip install -e .} 后，依序运行 \id{exposure}、\id{outcomes-range}（或 \id{build-all}）、\id{panel-did}、\id{did-regress --cluster twoway} 与 \id{did-regress-grid}。全量分析优先使用 Parquet。无本地 TeX 可用 Overleaf 编译。

\end{document}
